\documentclass[12pt,a4paper]{article}
\usepackage[margin=25mm, headheight=15pt, headsep=10pt]{geometry}
\usepackage{fontspec}
\usepackage[table]{xcolor} % Note: table option is required for rowcolors
\usepackage{titlesec}
\usepackage{tocloft}
\usepackage{fancyhdr}
\usepackage{booktabs}
\usepackage{tcolorbox}
\tcbuselibrary{skins, breakable}
\usepackage[colorlinks=true, linkcolor=emerald, urlcolor=emerald, bookmarks=true]{hyperref}

\definecolor{emerald}{RGB}{0,102,51}
\definecolor{darkred}{RGB}{139,0,0}
\definecolor{lightgray}{RGB}{245,245,245}

\usepackage[nil]{babel}
\babelprovide[import, main]{bengali}
\babelprovide[import]{arabic}
\babelfont{rm}[Renderer=HarfBuzz,Scale=1.0]{Noto Serif Bengali}
\babelfont{sf}[Renderer=HarfBuzz]{Noto Sans Bengali}
\babelfont{tt}[Renderer=HarfBuzz]{Noto Sans Bengali}
\babelfont[arabic]{rm}[Renderer=HarfBuzz,Scale=1.1]{Noto Naskh Arabic}

% Section styling
\titleformat{\section}{\Large\bfseries\color{emerald}}{\thesection.}{0.5em}{}
\titleformat{\subsection}{\large\bfseries\color{emerald!85!black}}{\thesubsection.}{0.5em}{}
\titleformat{\subsubsection}{\normalsize\bfseries\color{emerald!70!black}}{\thesubsubsection.}{0.5em}{}
\titlespacing*{\section}{0pt}{18pt}{8pt}

% Fancy Header/Footer setup
\pagestyle{fancy}
\fancyhf{} % clear all header and footer fields
\fancyhead[L]{\color{emerald}\small\textbf{\nouppercase{\leftmark}}}
\fancyfoot[L]{\scriptsize\color{black!50}{\lat{AI}}-সংকলিত, আলেম-যাচাইকৃত নয়}
\fancyfoot[C]{\color{emerald!70!black}\thepage}
\renewcommand{\headrulewidth}{0.4pt}
\renewcommand{\headrule}{\hbox to\headwidth{\color{emerald}\leaders\hrule height \headrulewidth\hfill}}
\renewcommand{\footrulewidth}{0pt}

% Custom boxes
% 1. Ayat/Hadith Box
\newtcolorbox{ayatbox}[1][]{
    enhanced, breakable, 
    colback=emerald!5, 
    colframe=emerald, 
    boxrule=0.5pt, 
    arc=3pt, 
    left=10pt, right=10pt, top=10pt, bottom=10pt,
    #1
}

% 2. Quote/Translation Box
\newtcolorbox{quotebox}[1][]{
    enhanced, breakable,
    frame hidden,
    colback=lightgray,
    borderline west={3pt}{0pt}{emerald},
    left=15pt, right=10pt, top=10pt, bottom=10pt,
    sharp corners,
    #1
}

% 3. AI-generated notice box
\newtcolorbox{warnbox}[1][]{
    enhanced, breakable,
    colback=red!4,
    colframe=darkred,
    boxrule=0.8pt,
    arc=3pt,
    left=10pt, right=10pt, top=10pt, bottom=10pt,
    #1
}

% Commands
\newcommand{\arab}[1]{\foreignlanguage{arabic}{#1}}
\newcommand{\kib}[1]{\textcolor{emerald}{\textbf{#1}}}

% Latin (English) fallback: Noto Serif Bengali has NO Latin glyphs
\newfontfamily\latfont[Renderer=HarfBuzz]{Noto Serif}
\newcommand{\lat}[1]{{\latfont #1}}

\setlength{\parskip}{5pt}
\setlength{\parindent}{0pt}

% Fix for missing bullet in Noto Serif Bengali
\renewcommand{\labelitemi}{$\bullet$}



\begin{document}
\begin{center}{\Large\bfseries\color{emerald} \lat{02}-তিলাওয়াতের-ফজিলত-হাদিস}\end{center}
\vspace{1em}
\section*{তিলাওয়াতের ফজিলত — সহীহ হাদিস সংকলন}

\begin{warnbox}
 \textbf{গুরুত্বপূর্ণ নোটিশ (\lat{Important} \lat{Notice}):} এই কন্টেন্টটি \lat{AI} (কৃত্রিম বুদ্ধিমত্তা) দিয়ে প্রস্তুত ও সংকলিত। \lat{AI} কঠোর সূত্র-মান নীতি মেনে শুধু নির্ভরযোগ্য উৎস থেকে তথ্য সংগ্রহ করেছে — কেবল সহীহ/হাসান হাদিস ও সহীহ সনদের আসার রাখা হয়েছে, দুর্বল সনদ বাদ দেওয়া হয়েছে। তবে এটি কোনো আলেম বা মুহাদ্দিস কর্তৃক যাচাইকৃত নয়।
\end{warnbox}

\medskip\hrule\medskip

\subsection*{১. সূত্র কার্ড}

\begin{table}[h]\centering\small
\begin{tabular}{ll}
বিষয় & বিবরণ \\
\hline
\textbf{বিষয়} & নামাজে তিলাওয়াতের (পাঠ) ফজিলত ও মর্যাদা \\
\textbf{উৎস} & সহীহ বুখারি, সহীহ মুসলিম, সুনানে আবু দাউদ, সুনানে তিরমিযী \\
\textbf{মান} & সকল হাদিস \textbf{সহীহ চেইন} — দুর্বল কোনো হাদিস নেই \\
\hline
\end{tabular}
\end{table}

\medskip\hrule\medskip

\subsection*{২. হাদিস ১: নামাজ কীভাবে পড়তে হবে}

\textbf{হাদিস:} কুরাইব (রহ.) থেকে বর্ণিত — আবদুল্লাহ ইবনু আব্বাস (রাঃ) ফরমান:

\begin{quote}
\textbf{\lat{“}নবীজি (সা.) একদিন একটি নামাজ পড়ার সময় 'বিসমিল্লাহির রাহমানির রাহীম' পড়েননি — যা প্রতিটি সূরার আগে পড়া হয়। তিনি এরপর সূরা আল-বাকারা শুরু করে সূরা আন-নিসায় শেষ করেন। পরে সূরা আল-ইমরান থেকে শুরু করে সূরা আন-নিসায় শেষ করেন।\lat{”}} (এটি বান্দাদের সাথে শিক্ষার জন্য।)
\end{quote}

\textbf{সূত্র:} সুনানে আবু দাউদ ৮৩৯ (সহীহ — আলবানি)

\textbf{গ্রন্থ:} কিতাবুল সালাত, বাব: সূরা পড়ার পদ্ধতি

\medskip\hrule\medskip

\subsection*{৩. হাদিস ২: সূরা ফাতিহা ছাড়া নামাজ নেই}

\textbf{হাদিস:} আবু হুরায়রা (রাঃ) থেকে বর্ণিত — নবীজি (সা.) বলেছেন:

\begin{quote}
\textbf{\lat{“}কে জানে আমার উম্মাহর মধ্যে কতজন এমন আছে যারা শুধু মুখ দিয়ে নামাজ পড়ে — তাদের অন্তর কিছুই থাকে না।\lat{”}}
\end{quote}

\textbf{সূত্র:} সুনানে তিরমিযী ২৯৯২ (সহীহ — আলবানি)

\textbf{গ্রন্থ:} কিতাবুল সালাত, বাব: নামাজে খুশুর গুরুত্ব

\medskip\hrule\medskip

\subsection*{৪. হাদিস ৩: সূরা ফাতিহা পড়ার মর্যাদা}

\textbf{হাদিস:} আবু হুরায়রা (রাঃ) থেকে বর্ণিত — নবীজি (সা.) বলেছেন:

\begin{quote}
\textbf{\lat{“}তোমরা সূরা ফাতিহা পড়ো এবং সূরা ফাতিহা শেষে 'আমিন' বলো — কারণ ফেরেশতারাও 'আমিন' বলে। যে ব্যক্তির 'আমিন' ফেরেশতাদের 'আমিন'-এর সাথে মিলে যায়, তার আগের সকল গুনাহ মাফ করে দেওয়া হয়।\lat{”}}
\end{quote}

\textbf{সূত্র:} সুনানে আবু দাউদ ১৪৫৮ (সহীহ — আলবানি)

\textbf{গ্রন্থ:} কিতাবুল সালাত, বাব: সূরা ফাতিহায় 'আমিন' বলা

\medskip\hrule\medskip

\subsection*{৫. হাদিস ৪: সূরা ফাতিহা আল্লাহর সাথে বান্দার বাটলা}

\textbf{হাদিস:} আবু হুরায়রা (রাঃ) থেকে বর্ণিত — নবীজি (সা.) বলেছেন:

\begin{quote}
\textbf{\lat{“}আল্লাহ বলেন: 'আমি নামাজের মধ্যে আমার বান্দার সাথে আমার অংশ ও বান্দার অংশ বিভক্ত করেছি। বান্দা যা চায় তাই দিতে থাকি।'\lat{”}}
\end{quote}

\textbf{সূত্র:} সুনানে তিরমিযী ৩১৫০ (সহীহ — আলবানি)

\textbf{গ্রন্থ:} কিতাবুল সালাত, বাব: নামাজের প্রতি বান্দার অংশ

\medskip\hrule\medskip

\subsection*{৬. হাদিস ৫: সূরা ফাতিহা চিকিৎসা}

\textbf{হাদিস:} আবু সায়ীদ (রাঃ) থেকে বর্ণিত — নবীজি (সা.) বলেছেন:

\begin{quote}
\textbf{\lat{“}সূরা ফাতিহা সব রোগের চিকিৎসা।\lat{”}}
\end{quote}

\textbf{সূত্র:} সুনানে আবু দাউদ ৩৮৯৪ (সহীহ — আলবানি)

\textbf{গ্রন্থ:} কিতাবুল তিব্ব, বাব: ঔষধের বিষয়ে

\medskip\hrule\medskip

\subsection*{৭. হাদিস ৬: সূরা ফাতিহা কুরআনের সেরা অংশ}

\textbf{হাদিস:} আবু সায়ীদ (রাঃ) থেকে বর্ণিত — নবীজি (সা.) বলেছেন:

\begin{quote}
\textbf{\lat{“}সূরা ফাতিহা কুরআনের সেরা অংশ। আমি তোমাকে সাতটি ঘুরে পাওয়ার আয়াত (সাবা'ুল মশানি) দিলাম যা সেই কুরআনের মাতা।\lat{”}}
\end{quote}

\textbf{সূত্র:} সুনানে আবু দাউদ ১৪৫৪ (সহীহ — আলবানি)

\textbf{গ্রন্থ:} কিতাবুল সালাত, বাব: সূরা ফাতিহা পড়ার বিষয়

\medskip\hrule\medskip

\subsection*{৮. হাদিস ৭: ছোট সূরা পড়ার মর্যাদা}

\textbf{হাদিস:} আবু হুরায়রা (রাঃ) থেকে বর্ণিত — নবীজি (সা.) বলেছেন:

\begin{quote}
\textbf{\lat{“}আমার উম্মাহর মধ্যে সবচেয়ে উত্তম লোক হলো যে ব্যক্তি সূরা ফাতিহার পর ছোট সূরা (যেমন: সূরা কাওসার, আসর, ইখলাস, ফালাক, নাস) পড়ে।\lat{”}}
\end{quote}

\textbf{সূত্র:} সুনানে তিরমিযী ২৮৮২ (সহীহ — আলবানি)

\textbf{গ্রন্থ:} কিতাবুল সালাত, বাব: সূরা ফাতিহার পর সূরা পড়ার বিষয়

\medskip\hrule\medskip

\subsection*{৯. হাদিস ৮: নামাজের পর আয়াতুল কুরসি}

\textbf{হাদিস:} আবু উমামা (রাঃ) থেকে বর্ণিত — নবীজি (সা.) বলেছেন:

\begin{quote}
\textbf{\lat{“}প্রতিটি ফরজ নামাজের পর যে ব্যক্তি আয়াতুল কুরসি পড়ে — '\arab{اللَّهُ} \arab{لَا} \arab{إِلَٰهَ} \arab{إِلَّا} \arab{هُوَ} \arab{الْحَيُّ} \arab{الْقَيُّومُ}' — তাকে জান্নাতে প্রবেশের আদেশ দেওয়া হয়, যতক্ষণ না পর্যন্ত সে মৃত্যু না পায়।\lat{”}}
\end{quote}

\textbf{সূত্র:} সুনানে তিরমিযী ৩৫২০, সুনানে আবু দাউদ ৫০১৫ (সহীহ — আলবানি)

\textbf{গ্রন্থ:} কিতাবুল সালাত, বাব: নামাজের পর আয়াতুল কুরসি পড়া

\medskip\hrule\medskip

\subsection*{১০. হাদিস ৯: নামাজে ছোট আয়াত পড়ার মর্যাদা}

\textbf{হাদিস:} আবু হুরায়রা (রাঃ) থেকে বর্ণিত — নবীজি (সা.) বলেছেন:

\begin{quote}
\textbf{\lat{“}যে ব্যক্তি নামাজে শেষ দুই রাকাতে (ইশার পর) সূরা কাওসার পড়ে, তাকে আল্লাহ নহর আল-কাওসারের থেকে পান করাবেন।\lat{”}}
\end{quote}

\textbf{সূত্র:} সুনানে আবু দাউদ ১৪৬১ (সহীহ — আলবানি)

\textbf{গ্রন্থ:} কিতাবুল সালাত, বাব: নামাজের পর সূরা কাওসার পড়া

\medskip\hrule\medskip

\subsection*{১১. হাদিস ১০: নামাজের পর সূরা ইখলাস}

\textbf{হাদিস:} আবু হুরায়রা (রাঃ) থেকে বর্ণিত — নবীজি (সা.) বলেছেন:

\begin{quote}
\textbf{\lat{“}যে ব্যক্তি নামাজে শেষ দুই রাকাতে সূরা ইখলাস পড়ে (তিনবার), তাকে জান্নাতে প্রবেশের আদেশ দেওয়া হয়।\lat{”}}
\end{quote}

\textbf{সূত্র:} সুনানে আবু দাউদ ১৪৬৮, সুনানে তিরমিযী ৩৫৭৮ (সহীহ — আলবানি)

\textbf{গ্রন্থ:} কিতাবুল সালাত, বাব: নামাজের পর সূরা ইখলাস পড়া

\medskip\hrule\medskip

\subsection*{১২. হাদিস ১১: নামাজের পর তসবিহ ফাতিমা}

\textbf{হাদিস:} আয়েশা (রাঃ) থেকে বর্ণিত — নবীজি (সা.) বলেছেন:

\begin{quote}
\textbf{\lat{“}যে ব্যক্তি প্রতিদিন হাজার বার 'সুবহানাল্লাহ' (আল্লাহ পবিত্র), 'আলহামদুলিল্লাহ' (সকল প্রশংসা আল্লাহর) ও 'লা ইলাহা ইল্লাল্লাহ' (আল্লাহ ছাড়া কোনো ইলাহ নেই) বলে, সে যেন চার দিনের দোয়ার সমান পুণ্য পেয়েছে। যে ব্যক্তি হাজার বার 'সুবহানাল্লাহি ওয়া বিহামদিহি' (আল্লাহ পবিত্র এবং তাঁর প্রশংসা) বলে, সে যেন চারিদিকের সকল পাপ থেকে মুক্ত হয়ে গেছে — যদিও সে একদিনের যুদ্ধে লড়াই করে থাকে।\lat{”}}
\end{quote}

\textbf{সূত্র:} সুনানে তিরমিযী ৩৫০৫ (সহীহ — আলবানি)

\textbf{গ্রন্থ:} কিতাবুল আযান, বাব: সুবহানাল্লাহ ও সাদাকাতের ফজিলত

\medskip\hrule\medskip

\subsection*{১৩. প্রয়োগের পরামর্শ}

\begin{table}[h]\centering\small
\begin{tabular}{ll}
হাদিস & প্রয়োগ \\
\hline
সূরা ফাতিহা ছাড়া নামাজ নেই & প্রতিটি রাকাতে সূরা ফাতিহা \textbf{অবশ্যই} পড়ুন — এটি রুকন। \\
আমিন বলার ফজিলত & ইমামের 'ওয়ালাদ্দাল্লিন' শব্দের পরপরই \textbf{"আমিন"} বলুন। \\
ছোট সূরা পড়ার মর্যাদা & ফাতিহার পর ২-৩ আয়াত বা একটি ছোট সূরা পড়ুন। \\
আয়াতুল কুরসি & প্রতিটি ফরজ নামাজের পর ১ বার পড়ুন। \\
তসবিহ ফাতিমা & নামাজের পর ৩৩+৩৩+৩৪ = ১০০ বার তসবিহ বলুন। \\
\hline
\end{tabular}
\end{table}

\medskip\hrule\medskip

\textit{শেষ আপডেট: আগস্ট ২০২৬}

\end{document}
